\documentclass[aspectratio=169,utf8,t]{beamer}

% Idioma e Codificacao
\usepackage[brazilian]{babel}
\usepackage{booktabs}
\usepackage{array}
\usepackage[normalem]{ulem}
\usepackage{tikz}
\usetikzlibrary{shapes,shapes.geometric,arrows,arrows.meta,positioning,calc,fit}

% Carrega o Tema Modelo dos Scouts (Widescreen 16:9, Nunito Sans)
\usetheme{scouts}

\definecolor{amber}{RGB}{245, 158, 11}

% Metadados Institucionais
\author{Prof. Me. Reginaldo Fernandes}
\institute{Instituto Federal do Ceará -- Campus Tauá \\ Tecnologia em Telemática}
\title{Banco de Dados}
\subtitle{Aula 10: Mapeamento Lógico Relacional (Parte 1) -- Entidades Fortes, Fracas e Relacionamentos 1:1 e 1:N}
\date{27 de Outubro de 2026}

\begin{document}

% ------------------------------------------------------------------------------
% Capa da Apresentacao
% ------------------------------------------------------------------------------
\changecolours[inverse]{ScoutGreen}
\frame[plain]{\titlepage}

% ------------------------------------------------------------------------------
% Sumario Executivo da Aula
% ------------------------------------------------------------------------------
\changecolours{ScoutGreen}
\begin{frame}{Agenda do Encontro}
  \begin{columns}[t]
    \begin{column}{0.48\textwidth}
      \scriptsize
      \tableofcontents[sections={1-5}]
    \end{column}
    \begin{column}{0.48\textwidth}
      \scriptsize
      \tableofcontents[sections={6-9}]
    \end{column}
  \end{columns}
\end{frame}

% ------------------------------------------------------------------------------
% Modulos Tematicos da Aula
% ------------------------------------------------------------------------------
\changecolours{ScoutGreen}
\section[Ponte e Objetivos]{Ponte Pedagógica e Objetivos de Aprendizagem}

% ------------------------------------------------------------------------------
\twocolframe[title={Do DER Conceitual ao Esquema Lógico Relacional},leftcol=7.0cm,rightcol=7.0cm]{%
  \alert{Onde Estávamos (Aulas 08 e 09)}\par
  \itemseps[topsepi=0.5mm,itemsepi=1mm,labelsep=2mm,leftmargini=3mm]
  \begin{itemize}
    \item \textbf{Aula 08 (Modelagem Conceitual CASE):} Construímos o diagrama DER/EER completo no Draw.io e brModelo, livre de detalhes de armazenamento.
    \item \textbf{Aula 09 (Modelo Relacional de Codd):} Dominamos a anatomia matemática das relações, domínios atômicos (1FN), chaves ($SK, CK, PK, AK, FK$) e as 4 restrições formais de integridade.
    \item \textbf{A Pergunta Central de Hoje:} Como converter sistematicamente cada figura geométrica do DER em relações matemáticas sem perder nenhuma regra de negócio?
  \end{itemize}
}{%
  \alert{Onde Chegaremos Hoje (Aula 10)}\par
  \itemseps[topsepi=0.5mm,itemsepi=1mm,labelsep=2mm,leftmargini=3mm]
  \begin{itemize}
    \item \textbf{Projeto Lógico Relacional:} Iniciamos a \textbf{Unidade IV (N1)} aplicando o algoritmo formal de transformação \textbf{MER $\longrightarrow$ Modelo Relacional}.
    \item \textbf{Escopo da Parte 1 (Passos 1 a 4):}\\
      -- \textbf{Passo 1:} Entidades Fortes e atributos (simples, compostos e derivados).\\
      -- \textbf{Passo 2:} Entidades Fracas e Relacionamentos Identificadores.\\
      -- \textbf{Passo 3:} Relacionamentos Binários $1:1$.\\
      -- \textbf{Passo 4:} Relacionamentos Binários $1:N$ e Auto-Relacionamentos.
  \end{itemize}
}

% ------------------------------------------------------------------------------
\twocolframe[title={Objetivos de Aprendizagem da Aula 10},leftcol=7.0cm,rightcol=7.0cm]{%
  \alert{Competências Teóricas e Analíticas}\par
  \itemseps[topsepi=0.5mm,itemsepi=1mm,labelsep=2mm,leftmargini=3mm]
  \begin{itemize}
    \item Compreender a lógica do \textbf{Algoritmo de Mapeamento de 7 Passos} de Elmasri \& Navathe e as regras de fusão/propagação de chaves de Carlos Alberto Heuser.
    \item Explicar por que atributos compostos devem ser achatados (\textit{flattened}) e atributos derivados omitidos das relações base para atender à \textbf{1ª Forma Normal (1FN)}.
    \item Justificar a escolha do lado de \textbf{Participação Total} ou de \textbf{Cardinalidade $N$} para posicionar Chaves Estrangeiras evitando valores \texttt{NULL}.
  \end{itemize}
}{%
  \alert{Habilidades Práticas de Modelagem}\par
  \itemseps[topsepi=0.5mm,itemsepi=1mm,labelsep=2mm,leftmargini=3mm]
  \begin{itemize}
    \item Transformar entidades regulares e fracas do DER em esquemas relacionais textuais e diagramas de esquema com arcos referenciais ($FK \longrightarrow PK$).
    \item Aplicar e comparar as \textbf{3 estratégias de mapeamento de relacionamentos $1:1$} (Chave Estrangeira, Fusão em Relação Única e Relação Própria).
    \item Mapear relacionamentos $1:N$ regulares e recursivos (auto-relacionamentos), renomeando chaves estrangeiras com base nos papéis (\textit{roles}).
  \end{itemize}
}

\changecolours{ScoutGreen}
\section[Algoritmo de 7 Passos]{Visão Geral do Algoritmo de Mapeamento MER $\to$ Relacional}

% ------------------------------------------------------------------------------
\begin{frame}{Pipeline de Transformação: Do Conceitual ao Lógico}
  \vspace{-0.15cm}
  \begin{center}
    \tikz[scale=0.68, transform shape]{
      % Box 1: MER/DER
      \node[draw=ScoutGreen, fill=ScoutGreen!15, very thick, rounded corners=4pt, minimum width=4.5cm, minimum height=2.2cm, align=center, font=\sffamily\scriptsize] (mer) at (-5.6, 0) {
        \textbf{\color{ScoutGreen}NÍVEL CONCEITUAL (MER/DER)}\\[3pt]
        $\bullet$ Retângulos (Entidades Fortes/Fracas)\\
        $\bullet$ Losangos ($1:1$, $1:N$, $N:M$)\\
        $\bullet$ Elipses (Compostos, Multival., Deriv.)\\
        $\bullet$ Hierarquias EER ($d/o$, Total/Parcial)
      };

      % Box 2: Algoritmo
      \node[draw=ScoutNavy, fill=ScoutNavy!12, very thick, rounded corners=4pt, minimum width=4.3cm, minimum height=2.2cm, align=center, font=\sffamily\scriptsize] (alg) at (0, 0) {
        \textbf{\color{ScoutNavy}ALGORITMO DE MAPEAMENTO}\\[3pt]
        \textbf{Elmasri \& Navathe (7 Passos)}\\
        \textbf{+ Heuser (Regras de Fusão)}\\[2pt]
        $\bullet$ \textbf{Aula 10:} Passos 1, 2, 3 e 4\\
        $\bullet$ \textbf{Aula 11:} Passos 5, 6, 7 e EER
      };

      % Box 3: Modelo Relacional
      \node[draw=ScoutGreen, fill=amber!15, very thick, rounded corners=4pt, minimum width=4.5cm, minimum height=2.2cm, align=center, font=\sffamily\scriptsize] (rel) at (5.6, 0) {
        \textbf{\color{ScoutNavy}NÍVEL LÓGICO RELACIONAL}\\[3pt]
        $\bullet$ Relações em 1FN (Atributos Atômicos)\\
        $\bullet$ Chaves Primárias (\underline{PK}) e Candidatas\\
        $\bullet$ Chaves Estrangeiras ($FK \longrightarrow PK$)\\
        $\bullet$ Restrições de \texttt{NULL} e Unicidade
      };

      \draw[-{Latex[length=2.8mm]}, ultra thick, ScoutGreen] (mer) -- (alg);
      \draw[-{Latex[length=2.8mm]}, ultra thick, ScoutNavy] (alg) -- (rel);
    }
  \end{center}
  \vspace{0.05cm}
  \begin{columns}[t]
    \begin{column}{0.48\textwidth}
      \footnotesize\alert{1. Preservação de Informação:}\par
      \scriptsize
      $\bullet$ Toda regra estrutural expressa no DER deve ter um equivalente exato no Esquema Relacional (relação, atributo atômico, chave primária, chave única ou chave estrangeira).
    \end{column}
    \begin{column}{0.48\textwidth}
      \footnotesize\alert{2. Minimização de Valores \texttt{NULL}:}\par
      \scriptsize
      $\bullet$ Ao decidir onde posicionar uma Chave Estrangeira ($FK$), sempre priorizamos a relação que possui \textbf{Participação Total $(1, 1)$} para evitar colunas repletas de valores nulos!
    \end{column}
  \end{columns}
\end{frame}

% ------------------------------------------------------------------------------
\begin{frame}{Panorama dos 7 Passos de Elmasri \& Navathe}
  \vspace{-0.18cm}
  \renewcommand{\arraystretch}{1.22}
  \resizebox{\textwidth}{!}{%
    \begin{tabular}{@{}ccp{4.8cm}p{6.8cm}c@{}}
      \toprule
      \textbf{Passo} & \textbf{Alvo no DER Conceitual} & \textbf{Construto Gráfico de Origem} & \textbf{Ação no Esquema Lógico Relacional} & \textbf{Aula} \\
      \midrule
      \textbf{Passo 1} & Entidades Regulares (Fortes) & Retângulo simples + Elipses simples, compostas e derivadas & Cria uma relação base; achata atributos compostos em subatributos atômicos e omite derivados. & \textbf{\color{ScoutGreen}Aula 10} \\
      \textbf{Passo 2} & Entidades Fracas ($\diamond\diamond$) & Retângulo duplo + Losango duplo + Discriminador tracejado & Cria relação para a fraca; importa a PK da forte como FK; define PK Composta ($\underline{FK + \text{Discriminador}}$). & \textbf{\color{ScoutGreen}Aula 10} \\
      \textbf{Passo 3} & Relacionamentos Binários $1:1$ & Losango $1:1$ (Parcial/Total ou Total/Total) & Propaga a PK para o lado de Participação Total (ou funde em relação única se ambos forem Totais). & \textbf{\color{ScoutGreen}Aula 10} \\
      \textbf{Passo 4} & Relacionamentos Binários $1:N$ & Losango $1:N$ (Regular ou Auto-Relacionamento Recursivo) & Propaga a PK do \textbf{lado $1$} como Chave Estrangeira ($FK$) para dentro da relação do \textbf{lado $N$}. & \textbf{\color{ScoutGreen}Aula 10} \\
      \midrule
      \textbf{Passo 5} & Relacionamentos Binários $N:M$ & Losango $N:M$ com ou sem atributos próprios & Cria \textbf{nova Relação Associativa} cuja PK é composta pelas FKs de ambos os lados. & \textit{Aula 11} \\
      \textbf{Passo 6} & Atributos Multivalorados & Elipse dupla ou pino $(0, n) / (1, n)$ & Cria \textbf{nova Relação Independente} contendo a FK da entidade proprietária + o valor do atributo. & \textit{Aula 11} \\
      \textbf{Passo 7} & Rel. $n$-ários e Herança EER & Losango ternário e Círculo/Triângulo EER & Cria relação associativa para grau $\ge 3$ e aplica uma das 4 opções de mapeamento EER. & \textit{Aula 11} \\
      \bottomrule
    \end{tabular}%
  }
\end{frame}

\changecolours{ScoutGreen}
\section[Passo 1: Entidades Fortes]{Passo 1: Mapeamento de Entidades Fortes e Seus Atributos}

% ------------------------------------------------------------------------------
\twocolframe[title={Passo 1: Regras para Entidades Fortes e Atributos},leftcol=7.0cm,rightcol=7.0cm]{%
  \alert{Como Mapear a Entidade Regular (Forte)}\par
  \itemseps[topsepi=0.5mm,itemsepi=1mm,labelsep=2mm,leftmargini=3mm]
  \begin{itemize}
    \item \textbf{Criação da Relação Base:} Para cada entidade regular (forte) $E$ do DER, cria-se um esquema de relação $R$ contendo todos os seus atributos simples.
    \item \textbf{Eleição da Chave Primária ($PK$):} Se a entidade possuir mais de um atributo identificador (chaves candidatas no DER, ex: \texttt{id\_paciente} e \texttt{cpf}), escolhe-se um como \textbf{Chave Primária (\underline{PK})} e documentam-se os demais como \textbf{Chaves Alternativas ($AK$ / Únicas)}.
  \end{itemize}
}{%
  \alert{Tratamento de Atributos Especiais no Passo 1}\par
  \itemseps[topsepi=0.5mm,itemsepi=1mm,labelsep=2mm,leftmargini=3mm]
  \begin{itemize}
    \item \textbf{Atributos Compostos (Achatamento / \textit{Flattening}):} Para respeitar a atomicidade da \textbf{1FN}, o nome do atributo tronco (ex: \texttt{endereco}) \textbf{não entra na relação}; incluímos apenas seus componentes atômicos folhas (\texttt{end\_rua}, \texttt{end\_num}, \texttt{end\_cep}).
    \item \textbf{Atributos Derivados (Omissão Física):} Atributos tracejados (ex: \texttt{idade}) \textbf{não viram colunas} na relação base para evitar redundância e anomalias de atualização.
    \item \textbf{Atributos Multivalorados:} \textbf{Ignorados no Passo 1} (serão tratados exclusivamente no Passo 6, na Aula 11!).
  \end{itemize}
}

% ------------------------------------------------------------------------------
\begin{frame}{Transformação Visual do Passo 1: Entidade Forte}
  \vspace{-0.2cm}
  \begin{center}
    \tikz[scale=0.66, transform shape]{
      % Lado Esquerdo: DER Conceitual
      \node[font=\scriptsize\bfseries\sffamily, ScoutGreen] at (-4.9, 2.05) {ANTES: Modelo Conceitual (DER)};
      \node[draw=ScoutGreen, fill=ScoutGreen!15, thick, minimum width=2.1cm, minimum height=0.65cm, font=\bfseries\sffamily\scriptsize] (pac) at (-5.0, 0.35) {PACIENTE};

      \node[draw=ScoutNavy, fill=ScoutNavy!12, ellipse, font=\tiny\sffamily\bfseries, inner sep=1.4pt] (pk) at (-7.4, 1.35) {\underline{id\_prontuario}};
      \node[draw=black!80, fill=white, ellipse, font=\tiny\sffamily, inner sep=1.4pt] (cpf) at (-7.5, 0.35) {cpf};
      \node[draw=black!80, fill=white, ellipse, font=\tiny\sffamily, inner sep=1.4pt] (dt) at (-7.3, -0.65) {data\_nasc};
      \node[draw=purple!80!black, fill=purple!10, dashed, ellipse, font=\tiny\sffamily, inner sep=1.4pt] (id) at (-5.0, -0.85) {idade};

      \node[draw=black!80, fill=amber!20, ellipse, font=\tiny\sffamily\bfseries, inner sep=1.4pt] (end) at (-2.8, 0.65) {endereco};
      \node[draw=black!80, fill=white, ellipse, font=\tiny\sffamily, inner sep=1.2pt] (rua) at (-1.4, 1.45) {rua};
      \node[draw=black!80, fill=white, ellipse, font=\tiny\sffamily, inner sep=1.2pt] (num) at (-1.2, 0.65) {numero};
      \node[draw=black!80, fill=white, ellipse, font=\tiny\sffamily, inner sep=1.2pt] (cep) at (-1.4, -0.15) {cep};

      \draw[semithick] (pac) -- (pk);
      \draw[semithick] (pac) -- (cpf);
      \draw[semithick] (pac) -- (dt);
      \draw[semithick, dashed, purple!80!black] (pac) -- (id);
      \draw[semithick] (pac) -- (end);
      \draw[semithick] (end) -- (rua);
      \draw[semithick] (end) -- (num);
      \draw[semithick] (end) -- (cep);

      % Seta de Transformacao
      \draw[-{Latex[length=3mm]}, ultra thick, ScoutNavy] (-0.3, 0.45) -- node[above, font=\tiny\bfseries\sffamily, ScoutNavy] {Passo 1} (0.9, 0.45);

      % Lado Direito: Esquema Relacional
      \node[font=\scriptsize\bfseries\sffamily, ScoutNavy] at (5.3, 2.05) {DEPOIS: Esquema de Relação (1FN)};
      \node[draw=ScoutNavy, fill=ScoutNavy!10, thick, minimum width=6.8cm, minimum height=0.70cm, font=\sffamily\scriptsize] (relpac) at (5.3, 0.95) {
        \textbf{PACIENTE}(\underline{id\_prontuario}, cpf, data\_nasc, \textbf{\color{ScoutGreen}end\_rua}, \textbf{\color{ScoutGreen}end\_num}, \textbf{\color{ScoutGreen}end\_cep})
      };

      \node[draw=amber!80!black, fill=amber!12, rounded corners=3pt, text width=6.8cm, inner sep=4pt, font=\tiny\sffamily] at (5.3, -0.35) {
        $\bullet$ \textbf{O que aconteceu com \texttt{endereco}?} Foi \textbf{achatado} nos 3 atributos atômicos (\texttt{end\_rua}, \texttt{end\_num}, \texttt{end\_cep}).\\
        $\bullet$ \textbf{O que aconteceu com \texttt{idade}?} Foi \textbf{omitido} da relação base (calculado a partir de \texttt{data\_nasc}).\\
        $\bullet$ \textbf{Chaves:} \underline{\texttt{id\_prontuario}} é a \textbf{PK} e \texttt{cpf} é Chave Alternativa (\textbf{AK}).
      };
    }
  \end{center}
  \vspace{-0.15cm}
  \renewcommand{\arraystretch}{1.28}
  \resizebox{\textwidth}{!}{%
    \begin{tabular}{@{}llllll@{}}
      \toprule
      \textbf{\underline{id\_prontuario} (PK)} & \textbf{cpf (AK)} & \textbf{data\_nasc} & \textbf{\color{ScoutGreen}end\_rua (Atômico)} & \textbf{\color{ScoutGreen}end\_numero (Atômico)} & \textbf{\color{ScoutGreen}end\_cep (Atômico)} \\
      \midrule
      \texttt{501} & \texttt{'123.456.789-00'} & \texttt{'1990-04-15'} & \texttt{'Av. Chermont Alves'} & \texttt{'410'} & \texttt{'63660-000'} \\
      \texttt{502} & \texttt{'987.654.321-11'} & \texttt{'1985-11-22'} & \texttt{'Rua Abigail Cidrão'} & \texttt{'85'} & \texttt{'63660-000'} \\
      \bottomrule
    \end{tabular}%
  }
\end{frame}

\changecolours{ScoutGreen}
\section[Passo 2: Entidades Fracas]{Passo 2: Mapeamento de Entidades Fracas e Relacionamentos Identificadores}

% ------------------------------------------------------------------------------
\twocolframe[title={Passo 2: Regras para Entidades Fracas ($W$) e Donas ($E$)},leftcol=7.0cm,rightcol=7.0cm]{%
  \alert{O Algoritmo Formal do Passo 2}\par
  \itemseps[topsepi=0.5mm,itemsepi=1mm,labelsep=2mm,leftmargini=3mm]
  \begin{itemize}
    \item \textbf{1. Criação da Relação Subordinada:} Para cada entidade fraca $W$ (retângulo duplo) pertencente a uma entidade forte proprietária $E$, cria-se uma relação $R_W$ com seus atributos simples.
    \item \textbf{2. Importação da PK Proprietária:} Inclui-se em $R_W$, como \textbf{Chave Estrangeira ($FK$)}, o atributo que é Chave Primária da relação proprietária $R_E$.
    \item \textbf{3. Chave Primária Composta:} A PK de $R_W$ é formada pela combinação da $FK$ importada + o Discriminador de $W$: $\underline{\text{PK}_{W}} = \underline{\text{FK}_{E} + \text{Discriminador}_{W}}$.
  \end{itemize}
}{%
  \alert{Consequências de Integridade no Passo 2}\par
  \itemseps[topsepi=0.5mm,itemsepi=1mm,labelsep=2mm,leftmargini=3mm]
  \begin{itemize}
    \item \textbf{Absorção do Losango Duplo:} O relacionamento identificador ($\diamond\diamond$) já fica 100\% representado pela $FK$ embutida na PK de $R_W$; \textbf{não se cria tabela extra para ele}!
    \item \textbf{Proibição Absoluta de \texttt{NULL}:} Como a $FK$ integra a Chave Primária de $R_W$, pela Integridade de Entidade ela \textbf{jamais aceita \texttt{NULL}}.
    \item \textbf{Exclusão em Cascata:} Se a tupla proprietária em $R_E$ for excluída, todas as tuplas dependentes em $R_W$ são removidas em \textbf{Cascata}.
  \end{itemize}
}

% ------------------------------------------------------------------------------
\begin{frame}{Transformação Visual do Passo 2: Entidade Fraca}
  \vspace{-0.18cm}
  \begin{center}
    \tikz[scale=0.67, transform shape]{
      % ANTES: DER
      \node[font=\scriptsize\bfseries\sffamily, ScoutGreen] at (-4.5, 1.90) {ANTES: Entidade Fraca no DER};
      \node[draw=ScoutGreen, fill=ScoutGreen!15, thick, minimum width=2.0cm, minimum height=0.6cm, font=\bfseries\sffamily\scriptsize] (pac) at (-6.8, 0.5) {PACIENTE};
      \node[draw=ScoutNavy, fill=ScoutNavy!12, ellipse, font=\tiny\sffamily\bfseries, inner sep=1.3pt] (pkp) at (-6.8, 1.45) {\underline{id\_prontuario}};

      \node[diamond, aspect=1.8, draw=ScoutNavy, double, double distance=1.3pt, fill=ScoutNavy!15, thick, inner sep=1pt, font=\bfseries\sffamily\tiny] (rel) at (-4.5, 0.5) {POSSUI\_DEP};

      \node[draw=ScoutGreen, double, double distance=1.4pt, fill=ScoutGreen!15, thick, minimum width=2.1cm, minimum height=0.6cm, font=\bfseries\sffamily\scriptsize] (dep) at (-2.0, 0.5) {DEPENDENTE};
      \node[draw=ScoutNavy, fill=amber!25, ellipse, font=\tiny\sffamily\bfseries, inner sep=1.3pt] (disc) at (-3.1, -0.55) {\dashuline{nome\_dep}};
      \node[draw=black!80, fill=white, ellipse, font=\tiny\sffamily, inner sep=1.3pt] (par) at (-1.1, -0.55) {parentesco};

      \draw[semithick] (pac) -- (pkp);
      \draw[semithick] (pac) -- node[above, font=\tiny\bfseries] {1} (rel);
      \draw[semithick] ($(rel.east)+(0,0.04)$) -- node[above, font=\tiny\bfseries] {N} ($(dep.west)+(0,0.04)$);
      \draw[semithick] ($(rel.east)+(0,-0.04)$) -- ($(dep.west)+(0,-0.04)$);
      \draw[semithick] (dep) -- (disc);
      \draw[semithick] (dep) -- (par);

      % Seta
      \draw[-{Latex[length=3mm]}, ultra thick, ScoutNavy] (-0.5, 0.5) -- node[above, font=\tiny\bfseries\sffamily, ScoutNavy] {Passo 2} (0.7, 0.5);

      % DEPOIS: Esquema Relacional
      \node[font=\scriptsize\bfseries\sffamily, ScoutNavy] at (4.8, 1.90) {DEPOIS: Diagrama de Esquema Relacional};
      \node[draw=ScoutGreen, fill=ScoutGreen!15, thick, minimum width=6.2cm, minimum height=0.62cm, font=\sffamily\scriptsize] (rpac) at (4.8, 1.15) {
        \textbf{PACIENTE}(\underline{id\_prontuario}, cpf, nome\_paciente)
      };

      \node[draw=ScoutNavy, fill=amber!15, thick, minimum width=6.8cm, minimum height=0.62cm, font=\sffamily\scriptsize] (rdep) at (4.8, -0.15) {
        \textbf{DEPENDENTE}(\underline{\textbf{\color{ScoutNavy}id\_prontuario\_fk}, \dashuline{nome\_dep}}, parentesco)
      };

      \draw[-{Latex[length=2.2mm]}, very thick, ScoutNavy] (3.4, -0.46) -- (3.4, -0.78) -- (1.1, -0.78) -- (1.1, 1.15) -- (1.7, 1.15);
      \node[font=\tiny\bfseries\sffamily, ScoutNavy] at (4.8, -0.98) {PK Composta = FK do Titular + Discriminador (Política: Exclusão em Cascata)};
    }
  \end{center}
  \vspace{-0.15cm}
  \renewcommand{\arraystretch}{1.24}
  \resizebox{\textwidth}{!}{%
    \begin{tabular}{@{}llll@{}}
      \toprule
      \textbf{\underline{\color{ScoutNavy}id\_prontuario\_fk} (PK/FK)} & \textbf{\underline{nome\_dep} (PK Parcial)} & \textbf{parentesco} & \textbf{Observação sobre Unicidade na Chave Composta} \\
      \midrule
      \texttt{501} & \texttt{'Lucas'} & \texttt{'Filho'} & Par $\langle 501, \text{'Lucas'} \rangle$ é único na relação \texttt{DEPENDENTE}. \\
      \texttt{501} & \texttt{'Marina'} & \texttt{'Cônjuge'} & Mesmo titular \texttt{501}, mas discriminador \texttt{'Marina'} distinto (Válido!). \\
      \texttt{502} & \texttt{'Lucas'} & \texttt{'Filho'} & O nome \texttt{'Lucas'} pode se repetir para outro paciente titular (\texttt{502} $\neq$ \texttt{501})! \\
      \bottomrule
    \end{tabular}%
  }
\end{frame}

\changecolours{ScoutGreen}
\section[Passo 3: Relacionamentos 1:1]{Passo 3: Mapeamento de Relacionamentos Binários $1:1$}

% ------------------------------------------------------------------------------
\begin{frame}{Passo 3: As 3 Estratégias Formais para Relacionamentos $1:1$}
  \vspace{-0.15cm}
  \footnotesize\alert{Como escolher a abordagem correta segundo a Participação (Total vs. Parcial):}\par
  \vspace{0.12cm}
  \renewcommand{\arraystretch}{1.42}
  \resizebox{\textwidth}{!}{%
    \begin{tabular}{@{}p{3.6cm}p{3.4cm}p{6.6cm}@{}}
      \toprule
      \textbf{Estratégia de Mapeamento} & \textbf{Quando Utilizar (Condição no DER)} & \textbf{Procedimento Formal no Esquema Lógico Relacional} \\
      \midrule
      \textbf{1. Abordagem da Chave Estrangeira} \newline \textit{(Mais comum na prática)} &
      Um lado tem \textbf{Participação Total} $(1, 1)$ e o outro tem \textbf{Participação Parcial} $(0, 1)$ (ou ambos parciais). &
      Escolhe-se a relação $S$ com \textbf{Participação Total $(1, 1)$} para receber a Chave Primária de $T$ como \textbf{Chave Estrangeira ($FK$)}, além dos atributos próprios do losango. A $FK$ também deve ser marcada como \textbf{Chave Alternativa ($AK$ / Única)} para garantir o $1:1$! \\
      \textbf{2. Abordagem da Fusão em Relação Única} &
      \textbf{Ambos os lados} possuem \textbf{Participação Total obrigatória} $(1, 1) \longleftrightarrow (1, 1)$. &
      Como as duas entidades têm exatamente a mesma quantidade de tuplas e não existem separadamente, fundem-se $S$ e $T$ em uma \textbf{única relação $R_{ST}$}, escolhendo uma das PKs como PK principal e a outra como $AK$. \\
      \textbf{3. Abordagem da Relação Própria de Referência Cruzada} &
      \textbf{Ambos os lados} têm \textbf{Participação Parcial} $(0, 1) \longleftrightarrow (0, 1)$ e pouquíssimas tuplas se relacionam. &
      Cria-se uma \textbf{terceira relação $R_{rel}$} contendo as duas chaves estrangeiras ($FK_S$ e $FK_T$, uma como PK e outra como AK) para evitar centenas de valores \texttt{NULL} nas tabelas principais. \\
      \bottomrule
    \end{tabular}%
  }
\end{frame}

% ------------------------------------------------------------------------------
\begin{frame}{Transformação Visual do Passo 3: Relacionamento $1:1$}
  \vspace{-0.2cm}
  \begin{center}
    \tikz[scale=0.66, transform shape]{
      % DER 1:1
      \node[font=\scriptsize\bfseries\sffamily, ScoutGreen] at (0, 1.95) {DER Conceitual: Todo Departamento DEVE ter 1 Diretor, mas nem todo Médico dirige um Departamento};
      \node[draw=ScoutGreen, fill=ScoutGreen!15, thick, minimum width=2.2cm, minimum height=0.6cm, font=\bfseries\sffamily\scriptsize] (med) at (-4.2, 0.9) {MEDICO};
      \node[draw=ScoutNavy, fill=ScoutNavy!12, ellipse, font=\tiny\sffamily\bfseries, inner sep=1.2pt] (pkm) at (-6.5, 0.9) {\underline{matricula}};

      \node[diamond, aspect=1.8, draw=ScoutNavy, fill=ScoutNavy!12, thick, inner sep=1.2pt, font=\bfseries\sffamily\tiny] (dir) at (0, 0.9) {DIRIGE};
      \node[draw=black!80, fill=white, ellipse, font=\tiny\sffamily, inner sep=1.2pt] (dt) at (0, -0.05) {dt\_inicio};

      \node[draw=ScoutGreen, fill=ScoutGreen!15, thick, minimum width=2.6cm, minimum height=0.6cm, font=\bfseries\sffamily\scriptsize] (dep) at (4.4, 0.9) {DEPTO\_CLINICO};
      \node[draw=ScoutNavy, fill=ScoutNavy!12, ellipse, font=\tiny\sffamily\bfseries, inner sep=1.2pt] (pkd) at (7.0, 0.9) {\underline{cod\_depto}};

      \draw[semithick] (med) -- (pkm);
      \draw[semithick] (dep) -- (pkd);
      \draw[semithick] (dir) -- (dt);
      \draw[semithick] (med) -- node[above, font=\scriptsize\bfseries, ScoutNavy] {1 \texttt{(0, 1)}} node[below, font=\tiny\sffamily] {Parcial} (dir);
      \draw[semithick] ($(dir.east)+(0,0.04)$) -- node[above, font=\scriptsize\bfseries, ScoutGreen] {1 \texttt{(1, 1)}} ($(dep.west)+(0,0.04)$);
      \draw[semithick] ($(dir.east)+(0,-0.04)$) -- node[below, font=\tiny\sffamily\bfseries, ScoutGreen] {Total!} ($(dep.west)+(0,-0.04)$);
    }
  \end{center}
  \vspace{-0.15cm}
  \begin{columns}[t]
    \begin{column}{0.48\textwidth}
      \footnotesize\alert{Escolha Errada (FK no Lado Parcial \texttt{MEDICO}):}\par
      \vspace{0.08cm}
      \renewcommand{\arraystretch}{1.2}
      \resizebox{\linewidth}{!}{%
        \begin{tabular}{@{}llll@{}}
          \toprule
          \textbf{\underline{matricula}} & \textbf{nome} & \textbf{\color{red!80!black}cod\_depto\_dir\_fk} & \textbf{\color{red!80!black}dt\_inicio} \\
          \midrule
          \texttt{1001} & \texttt{'Dra. Ana'} & \texttt{10} & \texttt{'2025-01-10'} \\
          \texttt{1002} & \texttt{'Dr. Bruno'} & \textbf{\color{red!80!black}NULL} & \textbf{\color{red!80!black}NULL} \\
          \texttt{1003} & \texttt{'Dra. Carla'} & \textbf{\color{red!80!black}NULL} & \textbf{\color{red!80!black}NULL} \\
          \texttt{1004} & \texttt{'Dr. Diego'} & \textbf{\color{red!80!black}NULL} & \textbf{\color{red!80!black}NULL} \\
          \bottomrule
        \end{tabular}%
      }
      \vspace{0.08cm}
      \scriptsize
      $\bullet$ \textbf{Problema:} Como 90\% dos médicos não são diretores, a tabela \texttt{MEDICO} fica poluída de células \texttt{NULL}!
    \end{column}
    \begin{column}{0.49\textwidth}
      \footnotesize\alert{Escolha Correta (FK no Lado Total \texttt{DEPTO\_CLINICO}):}\par
      \vspace{0.08cm}
      \renewcommand{\arraystretch}{1.2}
      \resizebox{\linewidth}{!}{%
        \begin{tabular}{@{}llll@{}}
          \toprule
          \textbf{\underline{cod\_depto}} & \textbf{nome\_depto} & \textbf{\color{ScoutGreen}mat\_diretor\_fk (AK)} & \textbf{\color{ScoutGreen}dt\_inicio\_dir} \\
          \midrule
          \texttt{10} & \texttt{'Cardiologia'} & \texttt{1001} (Não Nulo!) & \texttt{'2025-01-10'} \\
          \texttt{20} & \texttt{'Pediatria'} & \texttt{1005} (Não Nulo!) & \texttt{'2024-08-01'} \\
          \texttt{30} & \texttt{'Ortopedia'} & \texttt{1009} (Não Nulo!) & \texttt{'2026-02-15'} \\
          \bottomrule
        \end{tabular}%
      }
      \vspace{0.08cm}
      \scriptsize
      $\bullet$ \textbf{Solução Perfeita:} \textbf{Zero valores \texttt{NULL}}! E \texttt{mat\_diretor\_fk} recebe restrição de unicidade (\textbf{AK}) para impedir que um médico dirija 2 departamentos.
    \end{column}
  \end{columns}
\end{frame}

\changecolours{ScoutGreen}
\section[Passo 4: Relacionamentos 1:N]{Passo 4: Mapeamento de Relacionamentos Binários $1:N$ e Auto-Relacionamentos}

% ------------------------------------------------------------------------------
\twocolframe[title={Passo 4: Regras para Relacionamentos $1:N$ e Recursivos},leftcol=7.0cm,rightcol=7.0cm]{%
  \alert{A Regra Universal do Lado $N$}\par
  \itemseps[topsepi=0.5mm,itemsepi=1mm,labelsep=2mm,leftmargini=3mm]
  \begin{itemize}
    \item \textbf{Propagação do Lado $1$ para o Lado $N$:} Para cada relacionamento binário $1:N$, identifica-se a relação $S$ que representa a entidade no \textbf{lado $N$} (aquela que se associa a no máximo 1 ocorrência da outra).
    \item \textbf{Adição da Chave Estrangeira ($FK$):} Inclui-se na relação $S$ (lado $N$), como \textbf{Chave Estrangeira ($FK$)}, a Chave Primária da relação $T$ (lado $1$), além dos atributos simples do losango!
    \item \textbf{Por que NUNCA no lado $1$?} Colocar a FK no lado $1$ exigiria múltiplos valores na mesma célula (violando a \textbf{1FN}!).
  \end{itemize}
}{%
  \alert{Auto-Relacionamento $1:N$ e Obrigatoriedade}\par
  \itemseps[topsepi=0.5mm,itemsepi=1mm,labelsep=2mm,leftmargini=3mm]
  \begin{itemize}
    \item \textbf{Auto-Relacionamento Recursivo $1:N$:} Quando uma entidade se relaciona consigo mesma (ex: \texttt{MEDICO} preceptor $1 \longleftrightarrow N$ \texttt{MEDICO} residente), a $FK$ entra na \textbf{própria relação} renomeada pelo papel (\texttt{mat\_preceptor\_fk} $\longrightarrow$ \underline{\texttt{matricula}}).
    \item \textbf{Obrigatoriedade da FK no Lado $N$:}\\
      -- Lado $N$ com \textbf{Participação Total} $(1, 1)$: $FK$ \textbf{Não aceita \texttt{NULL}}.\\
      -- Lado $N$ com \textbf{Participação Parcial} $(0, 1)$: $FK$ \textbf{admite \texttt{NULL}}.
  \end{itemize}
}

% ------------------------------------------------------------------------------
\begin{frame}{Transformação Visual do Passo 4: Relacionamentos $1:N$}
  \vspace{-0.18cm}
  \begin{center}
    \tikz[scale=0.66, transform shape]{
      % ANTES: DER 1:N + Auto-relacionamento
      \node[font=\scriptsize\bfseries\sffamily, ScoutGreen] at (-4.3, 2.05) {ANTES: DER com Relacionamento $1:N$ e Recursivo $1:N$};

      \node[draw=ScoutGreen, fill=ScoutGreen!15, thick, minimum width=2.4cm, minimum height=0.6cm, font=\bfseries\sffamily\scriptsize] (dep) at (-6.8, 0.85) {DEPTO\_CLINICO};
      \node[diamond, aspect=1.7, draw=ScoutNavy, fill=ScoutNavy!12, thick, inner sep=1.2pt, font=\bfseries\sffamily\tiny] (lota) at (-4.1, 0.85) {LOTA};
      \node[draw=ScoutGreen, fill=ScoutGreen!15, thick, minimum width=2.1cm, minimum height=0.6cm, font=\bfseries\sffamily\scriptsize] (med) at (-1.5, 0.85) {MEDICO};

      \draw[semithick] (dep) -- node[above, font=\scriptsize\bfseries, ScoutNavy] {1} node[below, font=\tiny\sffamily] {(0, n)} (lota);
      \draw[semithick] (lota) -- node[above, font=\scriptsize\bfseries, ScoutGreen] {N} node[below, font=\tiny\sffamily] {(0, 1)} (med);

      % Auto-relacionamento SUPERVISIONA
      \node[diamond, aspect=1.8, draw=ScoutNavy, fill=amber!20, thick, inner sep=1pt, font=\bfseries\sffamily\tiny] (sup) at (-1.5, -0.75) {SUPERVISIONA};
      \draw[semithick] (med.south west) -- (-2.4, -0.75) node[midway, left, font=\tiny\sffamily] {preceptor \textbf{1} (0,n)} -- (sup.west);
      \draw[semithick] (med.south east) -- (-0.6, -0.75) node[midway, right, font=\tiny\sffamily] {\textbf{N} (0,1) residente} -- (sup.east);

      % Seta
      \draw[-{Latex[length=3mm]}, ultra thick, ScoutNavy] (0.5, 0.45) -- node[above, font=\tiny\bfseries\sffamily, ScoutNavy] {Passo 4} (1.6, 0.45);

      % DEPOIS: Diagrama de Esquema
      \node[font=\scriptsize\bfseries\sffamily, ScoutNavy] at (5.4, 2.05) {DEPOIS: Esquema Relacional (FKs no Lado $N$)};
      \node[draw=ScoutGreen, fill=ScoutGreen!15, thick, minimum width=6.4cm, minimum height=0.62cm, font=\sffamily\scriptsize] (rdep) at (5.4, 1.32) {
        \textbf{DEPTO\_CLINICO}(\underline{cod\_depto}, nome\_depto, bloco)
      };

      \node[draw=ScoutNavy, fill=ScoutNavy!12, thick, minimum width=7.2cm, minimum height=0.62cm, font=\sffamily\scriptsize] (rmed) at (5.4, -0.35) {
        \textbf{MEDICO}(\underline{matricula}, nome, \textbf{\color{ScoutNavy}cod\_depto\_fk}, \textbf{\color{ScoutGreen}mat\_preceptor\_fk})
      };

      % Arco 1: cod_depto_fk -> cod_depto
      \draw[-{Latex[length=2.1mm]}, very thick, ScoutNavy] (6.4, -0.04) -- (6.4, 0.58) -- (3.1, 0.58) -- (3.1, 1.01);
      \node[font=\tiny\bfseries\sffamily, ScoutNavy] at (4.85, 0.35) {FK do Lado 1 (\texttt{DEPTO}) no Lado N (\texttt{MEDICO})};

      % Arco 2: mat_preceptor_fk -> matricula
      \draw[-{Latex[length=2.1mm]}, very thick, ScoutGreen] (8.1, -0.66) -- (8.1, -1.08) -- (3.1, -1.08) -- (3.1, -0.66);
      \node[font=\tiny\bfseries\sffamily, ScoutGreen] at (5.6, -1.30) {FK Recursiva do Papel Lado 1 (\texttt{preceptor}) no Lado N (\texttt{residente})};
    }
  \end{center}
  \vspace{-0.18cm}
  \renewcommand{\arraystretch}{1.24}
  \resizebox{\textwidth}{!}{%
    \begin{tabular}{@{}lllll@{}}
      \toprule
      \textbf{\underline{matricula} (PK)} & \textbf{nome} & \textbf{\color{ScoutNavy}cod\_depto\_fk} & \textbf{\color{ScoutGreen}mat\_preceptor\_fk} & \textbf{Interpretação Semântica da Tupla no Lado $N$} \\
      \midrule
      \texttt{1001} & \texttt{'Dra. Ana'} & \texttt{10} (Cardiologia) & \texttt{NULL} & Médica sênior lotada no depto. \texttt{10}; não possui preceptor. \\
      \texttt{1002} & \texttt{'Dr. Bruno'} & \texttt{10} (Cardiologia) & \texttt{1001} (Dra. Ana) & Residente lotado no depto. \texttt{10} e supervisionado por \texttt{1001}. \\
      \texttt{1003} & \texttt{'Dra. Carla'} & \texttt{10} (Cardiologia) & \texttt{1001} (Dra. Ana) & Ambos (\texttt{1002} e \texttt{1003}) apontam para \texttt{1001} (materializa o $1:N$!). \\
      \bottomrule
    \end{tabular}%
  }
\end{frame}

\changecolours{ScoutGreen}
\section[Estudo de Caso Integrado]{Estudo de Caso Integrado e Tabela de Decisão (Passos 1 a 4)}

% ------------------------------------------------------------------------------
\begin{frame}{Esquema Integrado (Passos 1 a 4): Clínica ``VidaPlena''}
  \vspace{-0.15cm}
  \begin{center}
    \tikz[scale=0.68, transform shape]{
      % 1. DEPARTAMENTO_CLINICO (Passo 1 + Passo 3 FK Diretor)
      \node[draw=ScoutGreen, fill=ScoutGreen!15, thick, minimum width=7.5cm, minimum height=0.65cm, font=\sffamily\scriptsize] (depto) at (-4.3, 1.35) {
        \textbf{DEPTO\_CLINICO}(\underline{cod\_depto}, nome, \textbf{\color{ScoutGreen}mat\_diretor\_fk [AK]}, \textbf{\color{ScoutGreen}dt\_inicio})
      };

      % 2. MEDICO (Passo 1 + Passo 4 FK Lotacao + Passo 4 FK Preceptor)
      \node[draw=ScoutNavy, fill=ScoutNavy!12, thick, minimum width=7.5cm, minimum height=0.65cm, font=\sffamily\scriptsize] (med) at (4.5, 1.35) {
        \textbf{MEDICO}(\underline{matricula}, cpf [AK], nome, \textbf{\color{ScoutNavy}cod\_depto\_fk}, \textbf{\color{ScoutGreen}mat\_preceptor\_fk})
      };

      % 3. PACIENTE (Passo 1 com endereco achatado)
      \node[draw=ScoutGreen, fill=ScoutGreen!15, thick, minimum width=7.5cm, minimum height=0.65cm, font=\sffamily\scriptsize] (pac) at (-4.3, -0.75) {
        \textbf{PACIENTE}(\underline{id\_prontuario}, cpf [AK], nome, end\_rua, end\_num, end\_cep)
      };

      % 4. DEPENDENTE (Passo 2 Entidade Fraca)
      \node[draw=ScoutNavy, fill=amber!15, thick, minimum width=7.5cm, minimum height=0.65cm, font=\sffamily\scriptsize] (dep) at (4.5, -0.75) {
        \textbf{DEPENDENTE}(\underline{\textbf{\color{ScoutNavy}id\_prontuario\_fk}, \dashuline{nome\_dep}}, data\_nasc, parentesco)
      };

      % Arco Passo 3: mat_diretor_fk -> MEDICO.matricula
      \draw[-{Latex[length=2.1mm]}, very thick, ScoutGreen] (-2.9, 1.68) -- (-2.9, 2.08) -- (1.7, 2.08) -- (1.7, 1.68);
      \node[font=\tiny\bfseries\sffamily, ScoutGreen] at (-0.6, 2.28) {Passo 3 ($1:1$ DIRIGE -- FK no Lado Total)};

      % Arco Passo 4a: cod_depto_fk -> DEPTO_CLINICO.cod_depto
      \draw[-{Latex[length=2.1mm]}, very thick, ScoutNavy] (5.3, 1.02) -- (5.3, 0.55) -- (-6.0, 0.55) -- (-6.0, 1.02);
      \node[font=\tiny\bfseries\sffamily, ScoutNavy] at (-0.3, 0.34) {Passo 4 ($1:N$ LOTA -- FK no Lado $N$)};

      % Arco Passo 4b: mat_preceptor_fk -> MEDICO.matricula
      \draw[-{Latex[length=2.1mm]}, very thick, ScoutGreen] (7.4, 1.68) -- (7.4, 2.08) -- (2.1, 2.08) -- (2.1, 1.68);
      \node[font=\tiny\bfseries\sffamily, ScoutGreen] at (4.8, 2.28) {Passo 4 ($1:N$ Recursivo)};

      % Arco Passo 2: id_prontuario_fk -> PACIENTE.id_prontuario
      \draw[-{Latex[length=2.1mm]}, very thick, ScoutNavy] (2.4, -1.08) -- (2.4, -1.52) -- (-5.8, -1.52) -- (-5.8, -1.08);
      \node[font=\tiny\bfseries\sffamily, ScoutNavy] at (-1.7, -1.74) {Passo 2 (Entidade Fraca -- PK Composta = FK + Discriminador)};
    }
  \end{center}
  \vspace{-0.05cm}
  \scriptsize
  $\bullet$ \textbf{Observe a Economia de Tabelas:} Tivemos \textbf{4 entidades} e \textbf{4 losangos de relacionamento} ($\diamond\diamond\text{POSSUI\_DEP}$, $\diamond\text{DIRIGE}$, $\diamond\text{LOTA}$, $\diamond\text{SUPERVISIONA}$), mas o Esquema Lógico gerou \textbf{apenas 4 relações}! Todos os 4 relacionamentos foram absorvidos por propagação de Chaves Estrangeiras!
\end{frame}

% ------------------------------------------------------------------------------
\begin{frame}{Quadro Resumo de Decisão para os Passos 1 a 4}
  \vspace{-0.15cm}
  \renewcommand{\arraystretch}{1.32}
  \resizebox{\textwidth}{!}{%
    \begin{tabular}{@{}p{4.0cm}p{4.2cm}p{5.5cm}p{4.2cm}@{}}
      \toprule
      \textbf{Construto no DER} & \textbf{Gera Nova Relação?} & \textbf{Onde Posicionar a Chave Estrangeira ($FK$)?} & \textbf{Restrições na $FK$} \\
      \midrule
      \textbf{Passo 1: Entidade Forte} \newline (Simples, compostos e derivados) &
      \textbf{Sim} (1 relação base para cada entidade forte). &
      Não gera FK própria no Passo 1 (achata compostos em subatributos e omite derivados). &
      PK $\neq$ \texttt{NULL} e Única (Integridade de Entidade). \\
      \textbf{Passo 2: Entidade Fraca} \newline + Rel. Identificador ($\diamond\diamond$) &
      \textbf{Sim} (apenas para a entidade fraca; losango duplo absorvido). &
      Na relação da \textbf{Entidade Fraca}, importando a PK da entidade forte proprietária. &
      \textbf{NÃO aceita \texttt{NULL}} (integra a PK composta) + Cascata. \\
      \textbf{Passo 3: Rel. Binário $1:1$} \newline Parcial $(0,1) \times$ Total $(1,1)$ &
      \textbf{Não} (absorvido por FK, ou fusão se ambos forem $(1,1)$). &
      Na relação do lado de \textbf{Participação Total $(1, 1)$} (junto aos atributos do losango). &
      \textbf{NÃO aceita \texttt{NULL}} + Exige \textbf{Unicidade ($AK$)}. \\
      \textbf{Passo 4: Rel. Binário $1:N$} \newline (Regular ou Recursivo) &
      \textbf{Não} (exceto se o lado $N$ for opcional e raríssimo). &
      Sempre na relação do \textbf{Lado $N$} (que possui cardinalidade máxima $1$ no par de Heuser). &
      Aceita \texttt{NULL} se parcial $(0,1)$; \textbf{Não Nulo} se total $(1,1)$. \\
      \bottomrule
    \end{tabular}%
  }
\end{frame}

\changecolours{ScoutGreen}
\section[Conclusão e Exercícios]{Checklist de Validação e Exercícios de Fixação}

% ------------------------------------------------------------------------------
\twocolframe[title={Checklist de Validação do Mapeamento (Passos 1 a 4)},leftcol=7.0cm,rightcol=7.0cm]{%
  \alert{Auditoria Antes de Avançar para a Aula 11}\par
  \itemseps[topsepi=0.5mm,itemsepi=1mm,labelsep=2mm,leftmargini=3mm]
  \begin{itemize}
    \item[\textcolor{ScoutGreen}{$\boxtimes$}] Todos os atributos compostos foram substituídos pelos seus subatributos atômicos folhas (1FN)?
    \item[\textcolor{ScoutGreen}{$\boxtimes$}] Os atributos derivados foram removidos das relações base e registrados no Dicionário Lógico?
    \item[\textcolor{ScoutGreen}{$\boxtimes$}] Toda entidade fraca possui PK composta pela FK da entidade forte + o seu discriminador?
    \item[\textcolor{ScoutGreen}{$\boxtimes$}] Nos relacionamentos $1:1$, a FK foi colocada no lado de Participação Total $(1, 1)$ e marcada como $AK$?
    \item[\textcolor{ScoutGreen}{$\boxtimes$}] Nos relacionamentos $1:N$, a FK do lado $1$ foi posicionada dentro da relação do lado $N$?
  \end{itemize}
}{%
  \alert{Ponte Pedagógica: O Que Vem na Aula 11?}\par
  \itemseps[topsepi=0.5mm,itemsepi=1mm,labelsep=2mm,leftmargini=3mm]
  \begin{itemize}
    \item \textbf{E os relacionamentos $N:M$ e atributos multivalorados?} Eles \textbf{não podem} ser mapeados apenas adicionando uma coluna de FK em uma tabela existente!
    \item \textbf{Aula 11 (03/11/2026) -- Mapeamento Lógico Relacional (Parte 2):}\\
      -- \textbf{Passo 5:} Relacionamentos $N:M$ (Tabelas Associativas com PK Composta).\\
      -- \textbf{Passo 6:} Atributos Multivalorados (Relação Independente).\\
      -- \textbf{Passo 7:} Relacionamentos Ternários e Herança EER!
  \end{itemize}
}

% ------------------------------------------------------------------------------
\begin{frame}{Exercícios Práticos de Fixação -- Aula 10}
  \vspace{-0.18cm}
  \scriptsize
  \itemseps[topsepi=0.5mm,itemsepi=0.8mm,labelsep=2mm,leftmargini=4mm]
  \begin{enumerate}
    \item \textbf{Passo 1 (Compostos e Derivados):} A entidade forte \texttt{CLIENTE} possui \underline{\texttt{id\_cliente}}, \texttt{cpf} (único), \texttt{nome\_completo} (composto por \texttt{nome} e \texttt{sobrenome}), \texttt{dt\_cadastro} e \texttt{tempo\_fidelidade} (derivado). Escreva o esquema relacional resultante indicando PK e AK.
    \item \textbf{Passo 2 (Entidade Fraca):} A entidade forte \texttt{INTERNACAO} (\underline{\texttt{num\_int}}, \texttt{dt\_entrada}) identifica a entidade fraca \texttt{EVOLUCAO\_DIARIA} (discriminador \dashuline{\texttt{seq\_dia}}, \texttt{parecer}) pelo relacionamento identificador \texttt{REGISTRA}. Construa os esquemas relacionais e desenhe o arco de FK.
    \item \textbf{Passo 3 (Relacionamento $1:1$):} Considere \texttt{TECNICO} $(0, 1)$ --- $\langle\text{UTILIZA}\rangle$ (com atributo \texttt{km\_atual}) --- $(1, 1)$ \texttt{VIATURA}. Explique em qual relação a FK deve ser posicionada para evitar valores \texttt{NULL} e escreva os esquemas resultantes.
    \item \textbf{Passo 4 ($1:N$ Regular e Recursivo):} Um \texttt{POP\_REGIONAL} (\underline{\texttt{id\_pop}}, \texttt{cidade}) concentra vários \texttt{ROTEADORES} (\underline{\texttt{id\_router}}, \texttt{ip}, \texttt{modelo}) em $1:N$. Além disso, \texttt{ROTEADOR} possui auto-relacionamento $1:N$ \texttt{UPLINK} (1 \textit{gateway} $\to N$ \textit{downstream}). Escreva os esquemas relacionais.
    \item \textbf{Aplicação ao Projeto AVP1:} Aplique os Passos 1, 2, 3 e 4 sobre o DER que você desenhou na Aula 08.
  \end{enumerate}
\end{frame}

\changecolours{ScoutGreen}
\section[Referências]{Referências Bibliográficas}

% ------------------------------------------------------------------------------
\begin{frame}{Referências Bibliográficas Recomendadas}
  \scriptsize
  \begin{itemize}
    \item \textbf{ELMASRI, Ramez; NAVATHE, Shamkant B.} \textit{Sistemas de Banco de Dados}. 7. ed. São Paulo: Pearson Education do Brasil, 2018. \textbf{(Capítulo 9: Mapeamento do Modelo ER e EER para o Modelo Relacional -- Algoritmo de 7 Passos)}.
    \vspace{0.12cm}
    \item \textbf{HEUSER, Carlos Alberto.} \textit{Projeto de Banco de Dados}. 6. ed. Porto Alegre: Bookman, 2009. \textbf{(Capítulo 4: Projeto Lógico -- Transformação do Modelo Conceitual em Relacional)}.
    \vspace{0.12cm}
    \item \textbf{SILBERSCHATZ, Abraham; KORTH, Henry F.; SUDARSHAN, S.} \textit{Sistema de Banco de Dados}. 7. ed. Rio de Janeiro: Elsevier, 2020. \textbf{(Capítulo 6: Projeto de Bancos de Dados e o Modelo E-R -- Redução para Esquemas Relacionais)}.
    \vspace{0.12cm}
    \item \textbf{CODD, Edgar Frank.} \textit{A Relational Model of Data for Large Shared Data Banks}. Communications of the ACM (CACM), New York, v. 13, n. 6, p. 377--387, June 1970.
    \vspace{0.12cm}
    \item \textbf{CHEN, Peter Pin-Shan.} \textit{The Entity-Relationship Model: Toward a Unified View of Data}. ACM Transactions on Database Systems (TODS), New York, v. 1, n. 1, p. 9--36, Mar. 1976.
  \end{itemize}
\end{frame}


\end{document}
